\section{Introduction}
\label{sec:intro}

\IEEEPARstart{C}{ardiovascular} diseases remain the leading cause of death worldwide, and the
Global Burden of Disease study estimated that their prevalence has almost doubled between 1990
and 2019 \cite{roth2020gbd}. The 12-lead electrocardiogram (ECG) is the first-line, low-cost and
non-invasive test for most of these conditions: it is recorded in primary care, emergency
departments and wards, and a single ten-second recording can reveal myocardial infarction (MI),
ischaemic ST/T changes (STTC), conduction disturbances (CD) and ventricular hypertrophy (HYP).
Because a large fraction of ECGs is acquired where no cardiologist is present, reliable automatic
interpretation directly affects how quickly a patient with an acute finding is identified and
referred.

\textit{Problem definition.} In current practice the ECG is interpreted visually by a physician,
often supported by rule-based software built into the electrocardiograph. Visual reading is
subjective and depends on training and workload, and the built-in algorithms are known to produce
frequent errors, especially for non-sinus rhythms and subtle morphological changes, so their
statements have to be over-read by an expert \cite{schlapfer2017}. A single ECG can, moreover,
carry several co-existing diagnoses, which makes the task inherently multi-label rather than a
choice between mutually exclusive classes. In this work the problem is therefore formulated as
the multi-label prediction of the five diagnostic superclasses NORM, MI, STTC, CD and HYP from the
raw 12-lead signal.

Deep learning has substantially changed automatic ECG analysis. Convolutional and recurrent
networks have reached cardiologist-level performance on single-lead ambulatory recordings
\cite{hannun2019} and on large 12-lead cohorts \cite{ribeiro2020}, and the public PTB-XL dataset
\cite{wagner2020ptbxl} has enabled reproducible benchmarks in which ResNet, Inception and LSTM
models reach macro areas under the ROC curve (AUROC) between 0.921 and 0.930 on the
superdiagnostic task \cite{strodthoff2021}. Subsequent work has added self-supervised pretraining
\cite{mehari2022ssl,na2024stmem}, state space models and patient metadata \cite{mehari2023s4}, and
ECG--text multimodal learning \cite{li2023mets,liu2024merl}. Three gaps nevertheless remain.
First, almost all of these models are trained on centrally pooled data, whereas in practice ECG
archives are held by separate hospitals that cannot share patient records. Second, the federated
learning (FL) methods that address this constraint \cite{rieke2020future,kairouz2021advances}, such as FedAvg \cite{mcmahan2017fedavg},
FedProx \cite{li2020fedprox} and SCAFFOLD \cite{karimireddy2020scaffold}, average model parameters
and therefore require every site to train the same architecture over many communication rounds.
Third, the clinical use of such models requires explanations and an explicit statement of when the
model should not decide on its own, which most ECG studies do not provide \cite{wagner2024xai}.

\textit{Motivation and novelty.} This paper studies a setting in which several institutions each
own a disjoint, patient-separated share of the training data and are free to choose their own
network architecture. Instead of averaging parameters, each institution trains its model once,
calibrates it, and contributes only its predicted class probabilities to a central gate, which
combines them as $P(y\mid x)=\sum_i w_i P_i(y\mid x)$ with weights on the probability simplex. This
one-shot, prediction-level form of federation needs a single communication step per query, places
no constraint on the client architectures, and never moves raw ECGs or parameters between sites. The
question addressed is deliberately conservative: whether the aggregated model is better than the
best individual client, and how close it comes to a model trained on all data pooled together. The
weighting scheme itself is selected on validation data with a rule that prefers the simplest
combiner unless a more complex one is clearly better, and the final model is evaluated once on the
held-out PTB-XL test fold over three random seeds.

The major contributions of this work are as follows.
\begin{enumerate}
  \item \textbf{One-shot federated aggregation of heterogeneous clients:} Four clients train
  structurally different networks (a 1D CNN, a 1D ResNet, a BiLSTM and a CNN-LSTM) on disjoint,
  patient-level shares of PTB-XL, and only calibrated probabilities are exchanged, so no parameter
  averaging and no common architecture are required.
  \item \textbf{Calibrated simplex gate with parsimony-first selection:} Equal weights, learned
  static weights and a per-patient dynamic gate are fitted on patient-grouped cross-validation
  predictions, and a more complex combiner is adopted only if it improves the validation macro AUROC
  by at least 0.002.
  \item \textbf{Rigorous multi-seed evaluation against strong references:} Over three seeds the
  selected gate reaches a test macro AUROC of 0.905 and improves on the best single client in every
  seed (paired bootstrap), and it is compared with traditional classifiers on 172 hand-crafted
  features and with FedAvg, FedProx, SCAFFOLD and FedNova on the same partitions, remaining within
  0.002 AUROC of FedAvg with a single communication round.
  \item \textbf{Client-level ablation and attribution:} All 15 client subsets are evaluated and
  the contribution of each institution is quantified with Shapley values, showing why learned
  weights that discard weaker clients lose to equal weighting.
  \item \textbf{Explainability and safe automation:} Integrated Gradients attribute each superclass
  prediction to leads and time, and per-finding rule-in/rule-out bands fitted on validation data
  automate 53.7\% of findings with a precision of 0.937 and an accuracy of 0.967 while referring the
  remaining findings to a cardiologist.
\end{enumerate}

The rest of this paper is organized as follows. Section~\ref{sec:related} reviews recent unimodal,
multimodal, federated and explainable approaches to ECG and clinical classification.
Section~\ref{sec:method} describes the proposed model, including preprocessing, the client
networks, the federated aggregation and the explainability methods. Section~\ref{sec:results}
presents the dataset, the experimental setup, the comparisons, the ablation experiments, the
federated and class-wise analyses and the explainability results. Section~\ref{sec:conclusion}
concludes the paper and outlines future work.
